\section{Intrinsic per-particle moments}\label{sec:intrinsic}
The conserved family has a nonlinear consequence without enlarging the
unbounded-test class. Let
\[
 M_\eps=\pi_0^\eps(1)=N/\mu,\qquad
 p_\eps=\pi_t^\eps(v)/M_\eps,\qquad
 x_\eps(t)=\pi_t^\eps(x)/M_\eps.
\]
On $N>0$ define the intensive intrinsic moments
\begin{equation}\label{eq:intrinsic-def}
 \begin{split}
 U_\eps&=\pi_t^\eps(|v|^2)/M_\eps-|p_\eps|^2,\\
 D_\eps(t)&=\pi_t^\eps(x\cdot v)/M_\eps-x_\eps(t)\cdot p_\eps,\\
 S_\eps(t)&=\pi_t^\eps(|x|^2)/M_\eps-|x_\eps(t)|^2.
 \end{split}
\end{equation}
Assign any fixed values on $N=0$. The mass coordinate of
Theorem~\ref{thm:main} implies $M_\eps\to1$ in probability, so that
exceptional probability tends to zero. Write
\[
 p_0=f_0(v),\quad x_0=f_0(x),\quad
 u_0=f_0(|v-p_0|^2),\quad
 d_0=f_0((x-x_0)\cdot(v-p_0)),\quad
 s_0=f_0(|x-x_0|^2).
\]
Their Boltzmann values are $(u_0,d_0+tu_0,s_0+2td_0+t^2u_0)$.
Define three effective initial tests
\begin{equation}\label{eq:effective-tests}
 \phi_U=|v-p_0|^2-u_0,\qquad
 \phi_D=(x-x_0)\cdot(v-p_0)-d_0,\qquad
 \phi_S=|x-x_0|^2-s_0.
\end{equation}
Each belongs to $\cQ$ and integrates to zero against $f_0$.

\begin{corollary}[Thermal energy and intrinsic spread]\label{cor:intrinsic}
For either flow, the Boltzmann-centered vector of
\eqref{eq:intrinsic-def}, scaled by $\eps^{-1}$, converges in
$D([0,T],\R^3)$, $J_1$, to
\begin{equation}\label{eq:intrinsic-limit}
 A(t)(Z^{\rm int}+a^{\rm int})+c(t),
 \qquad
 \Sigma^{\rm int}_{ij}=\int f_0\phi_i\phi_j,\qquad
 a_i^{\rm int}=a(\phi_i),
\end{equation}
where $Z^{\rm int}$ is centered Gaussian with covariance
$\Sigma^{\rm int}$, and $A,c$ are from \eqref{eq:three-limit}.
This convergence is joint with the exactly centered Schwartz field, with
$Z_i^{\rm int}=G_0(\eta\phi_i)$. It describes the weak limit and its
covariance; it does not add a true prelimit covariance assertion.
\end{corollary}
\begin{proof}
Mass and momentum are constant, and
$x_\eps(t)=x_\eps(0)+tp_\eps$.
Substitution in \eqref{eq:energy}--\eqref{eq:spread} gives, exactly,
\begin{align*}
 U_\eps(t)&=U_\eps(0),\\
 D_\eps(t)&=D_\eps(0)+tU_\eps(0)+\eps\cH_\eps(t)/M_\eps,\\
 S_\eps(t)&=S_\eps(0)+2tD_\eps(0)+t^2U_\eps(0)
                         +2\eps\int_0^t\cH_\eps(s)\dd s/M_\eps.
\end{align*}
Since $M_\eps\to1$, Proposition~\ref{prop:LLN} gives the same uniform
deterministic event term after division by $M_\eps$.

For clarity, denote the raw initial moments by
$(M,P,X,E,V,Q)$. The three rational maps are
$E/M-|P|^2/M^2$, $V/M-X\cdot P/M^2$, and
$Q/M-|X|^2/M^2$. At the deterministic point
$(1,p_0,x_0,f_0(|v|^2),f_0(x\cdot v),f_0(|x|^2))$
their differentials, represented as tests of an empirical perturbation, are
exactly \eqref{eq:effective-tests}. For example the first differential is
\[
 \delta E-2p_0\cdot\delta P
       +(2|p_0|^2-f_0(|v|^2))\delta M
 =\delta\pi\bigl(|v-p_0|^2-u_0\bigr).
\]
The same calculation gives the other two tests. The initial thirteen-vector
has an $O_{\Prob}(\eps)$ deviation and its scaled limit is its Gaussian
vector plus the vacancy bias. On $M$ bounded away from zero the rational
maps have bounded second derivatives in a fixed neighborhood. Thus their
$\eps^{-1}$-scaled Taylor remainders are $O_{\Prob}(\eps)\to0$.
The finite-dimensional delta method gives the initial limit
$Z^{\rm int}+a^{\rm int}$, jointly with the field. The exact triangular
identities and the uniform event LLN prove \eqref{eq:intrinsic-limit}.
\end{proof}

The observables are unchanged by a constant spatial translation or Galilean
boost. Normalization by the particle count removes the leading mass direction;
this explains the zero means of the effective tests under $f_0$, while the
initial vacancy functional $a(\phi_i)$ can still be nonzero. The corollary
does not require a separate positive-time CLT for these rational observables.
